\section{The Field-Mismatch Tax and When a Precompile Pays}
We formalize the \emph{field-mismatch tax} as the performance gap that arises when a cryptographic primitive's natural algebraic domain (e.g., arithmetic over a large prime field $\mathbb{F}_q$ or structured rings) does not coincide with the proof system's \emph{native} arithmetic, and derive from it a criterion for when a custom precompile lowers proving cost. We model a proof backend by its \emph{arithmetization regime}: (i) a native field $\mathbb{F}_p$, (ii) an encoding of the computation domain into $\mathbb{F}_p$ (possibly via multi-limb representations), and (iii) a constraint/execution representation (e.g., R1CS for static circuits or AIR/trace constraints for STARK-based zkVMs). This viewpoint separates two qualitatively different regimes.

In the \emph{field-aligned} regime typical of application-specific static circuits, the designer can select (or align to) a proof system whose scalar field matches the primitive's domain, so arithmetic operations are expressed directly: multiplications correspond to constant-cost constraints, and total constraint growth tracks the primitive's multiplicative complexity. In the \emph{mismatch} regime typical of general-purpose zkVMs, the native field $\mathbb{F}_{p_{\text{VM}}}$ is fixed by the VM/STARK design and is often much smaller than $\mathbb{F}_q$. Elements of $\mathbb{F}_q$ must then be represented as $k=\lceil \log_2(q)/w\rceil$ $w$-bit limbs, requiring carry propagation and modular reduction to emulate non-native arithmetic. We show that even a single non-native modular multiplication is lower-bounded by $\Omega(k)$ native multiplications (with schoolbook emulation realizing all $k^2$ cross-products, plus carries and reduction), yielding an inherent field-mismatch tax that cannot be eliminated by constant-factor engineering alone. The lower bound is a counting argument: in a limb-aligned routine the product depends on every one of the $k$ input limbs, while each native multiplication carries at most one of them, so at least $k$ native multiplications are forced.

For STARK-based zkVMs, this cost further amplifies at the proof-system level: limb-based arithmetic increases the executed instruction count and inflates the AIR execution trace length by $\Theta(k^2)$ under the schoolbook emulation the evaluated zkVMs use. Since zkVM prover time is dominated by low-degree testing and commitments and scales roughly as $O(W\cdot T\log T)$ in trace width $W$ and length $T$, arithmetic mismatch translates into a fundamental prover slowdown relative to field-aligned static circuits. This yields the thesis's criterion: a precompile lowers proving cost only on the field-mismatch axis, and it pays only when it removes more emulated trace area than the dedicated chip itself adds. A primitive already served by an existing precompile, or one native to the prover field, therefore does not warrant a bespoke chip even though one could be built.
